\section{Additional result tables}
\label{app:extra}

These tables complete Section~\ref{sec:results}; they are copied from the same analysis reports.

\begin{table}[H]
\centering
\caption{Headline runs of the main analysis: all features or one limb, raw or level-calibrated, logistic regression or forest,
with macro-F1, permutation $p$ and the 95th percentile of the permutation null.}
\label{tab:headline}
\begin{adjustbox}{max width=\textwidth}
{\scriptsize
\begin{tabular}{lllrrrrrrrr}
\toprule
\textbf{task} & \textbf{features} & \textbf{norm} & \textbf{model} & \textbf{n feat} & \textbf{n rows} & \textbf{bal. acc} & \textbf{macro F1} & \textbf{chance} & \textbf{perm p} & \textbf{null 95\%} \\
\midrule
height & all & calibrated & logreg & 110 & 260 & 53.7\% & 0.54 & 33\% & $<$ 0.001 & 37.8\% \\
height & all & raw & logreg & 110 & 260 & 53.3\% & 0.53 & 33\% & - & - \\
height & all & calibrated & forest & 110 & 260 & 58.9\% & 0.59 & 33\% & - & - \\
height & all & raw & forest & 110 & 260 & 49.4\% & 0.48 & 33\% & - & - \\
height & lead only & calibrated & logreg & 57 & 260 & 62.1\% & 0.62 & 33\% & $<$ 0.001 & 37.9\% \\
height & lead only & raw & logreg & 57 & 260 & 56.4\% & 0.56 & 33\% & - & - \\
height & lead only & calibrated & forest & 57 & 260 & 59.3\% & 0.59 & 33\% & - & - \\
height & lead only & raw & forest & 57 & 260 & 51.6\% & 0.50 & 33\% & - & - \\
height & trail only & calibrated & logreg & 52 & 260 & 49.1\% & 0.48 & 33\% & $<$ 0.001 & 38.0\% \\
height & trail only & raw & logreg & 52 & 260 & 53.6\% & 0.52 & 33\% & - & - \\
height & trail only & calibrated & forest & 52 & 260 & 59.2\% & 0.59 & 33\% & - & - \\
height & trail only & raw & forest & 52 & 260 & 43.9\% & 0.43 & 33\% & - & - \\
width & all & calibrated & logreg & 110 & 260 & 57.8\% & 0.58 & 50\% & 0.003 & 54.5\% \\
width & all & raw & logreg & 110 & 260 & 56.5\% & 0.56 & 50\% & - & - \\
width & all & calibrated & forest & 110 & 260 & 60.3\% & 0.60 & 50\% & - & - \\
width & all & raw & forest & 110 & 260 & 54.2\% & 0.54 & 50\% & - & - \\
width & lead only & calibrated & logreg & 57 & 260 & 54.1\% & 0.54 & 50\% & 0.099 & 55.0\% \\
width & lead only & raw & logreg & 57 & 260 & 52.9\% & 0.53 & 50\% & - & - \\
width & lead only & calibrated & forest & 57 & 260 & 61.0\% & 0.61 & 50\% & - & - \\
width & lead only & raw & forest & 57 & 260 & 56.1\% & 0.56 & 50\% & - & - \\
width & trail only & calibrated & logreg & 52 & 260 & 57.0\% & 0.57 & 50\% & 0.013 & 55.2\% \\
width & trail only & raw & logreg & 52 & 260 & 55.6\% & 0.55 & 50\% & - & - \\
width & trail only & calibrated & forest & 52 & 260 & 54.8\% & 0.55 & 50\% & - & - \\
width & trail only & raw & forest & 52 & 260 & 51.1\% & 0.51 & 50\% & - & - \\
\bottomrule
\end{tabular}
}
\end{adjustbox}

\end{table}

\begin{table}[H]
\centering
\caption{Height and width per segment and limb with the generic (level-calibrated) logistic regression.}
\label{tab:segmentslevel}
\begin{adjustbox}{max width=\textwidth}
{\scriptsize
\begin{tabular}{llllrrrrrrr}
\toprule
\textbf{task} & \textbf{features} & \textbf{norm} & \textbf{model} & \textbf{n feat} & \textbf{n rows} & \textbf{bal. acc} & \textbf{macro F1} & \textbf{chance} & \textbf{perm p} & \textbf{null 95\%} \\
\midrule
height & thigh & calibrated & logreg & 24 & 260 & 58.3\% & 0.58 & 33\% & 0.003 & 37.8\% \\
height & thigh (lead) & calibrated & logreg & 12 & 260 & 53.7\% & 0.53 & 33\% & 0.003 & 37.2\% \\
height & thigh (trail) & calibrated & logreg & 12 & 260 & 41.0\% & 0.39 & 33\% & 0.007 & 37.9\% \\
height & shank & calibrated & logreg & 24 & 260 & 52.3\% & 0.51 & 33\% & 0.003 & 37.7\% \\
height & shank (lead) & calibrated & logreg & 12 & 260 & 43.3\% & 0.43 & 33\% & 0.003 & 37.7\% \\
height & shank (trail) & calibrated & logreg & 12 & 260 & 51.2\% & 0.51 & 33\% & 0.003 & 38.4\% \\
height & foot angle & calibrated & logreg & 26 & 260 & 61.2\% & 0.60 & 33\% & 0.003 & 37.0\% \\
height & foot angle (lead) & calibrated & logreg & 13 & 260 & 66.5\% & 0.65 & 33\% & 0.003 & 37.9\% \\
height & foot angle (trail) & calibrated & logreg & 13 & 260 & 53.3\% & 0.52 & 33\% & 0.003 & 37.2\% \\
height & foot acc & calibrated & logreg & 6 & 260 & 44.9\% & 0.45 & 33\% & 0.003 & 37.9\% \\
height & foot acc (lead) & calibrated & logreg & 3 & 260 & 38.2\% & 0.36 & 33\% & 0.037 & 37.1\% \\
height & foot acc (trail) & calibrated & logreg & 3 & 260 & 47.0\% & 0.46 & 33\% & 0.003 & 37.4\% \\
height & ZHC trajectory & calibrated & logreg & 22 & 260 & 56.3\% & 0.57 & 33\% & 0.003 & 38.6\% \\
height & ZHC trajectory (lead) & calibrated & logreg & 11 & 260 & 60.0\% & 0.60 & 33\% & 0.003 & 37.6\% \\
height & ZHC trajectory (trail) & calibrated & logreg & 11 & 260 & 57.9\% & 0.58 & 33\% & 0.003 & 37.9\% \\
height & spatiotemporal & calibrated & logreg & 10 & 260 & 53.7\% & 0.53 & 33\% & 0.003 & 38.2\% \\
height & spatiotemporal (lead) & calibrated & logreg & 7 & 260 & 44.3\% & 0.44 & 33\% & 0.003 & 38.0\% \\
height & spatiotemporal (trail) & calibrated & logreg & 2 & 260 & 54.5\% & 0.54 & 33\% & 0.003 & 37.2\% \\
\bottomrule
\end{tabular}
}
\end{adjustbox}

\end{table}

\begin{table}[H]
\centering
\caption{Permutation tests of the separate-limb models (labels shuffled within participant).}
\label{tab:limbsperm}
\begin{adjustbox}{max width=\textwidth}
{\footnotesize
\begin{tabular}{lllrrr}
\toprule
\textbf{target} & \textbf{norm} & \textbf{limb} & \textbf{bal. acc} & \textbf{perm p} & \textbf{null 95th pct} \\
\midrule
height & level & lead & 62.1\% & 0.003 & 37.5\% \\
height & level & fusion & 56.1\% & 0.003 & 37.6\% \\
height & self & lead & 79.6\% & 0.003 & 38.5\% \\
height & self & fusion & 80.8\% & 0.003 & 38.8\% \\
width & level & lead & 54.1\% & 0.103 & 55.1\% \\
width & level & fusion & 56.1\% & 0.020 & 54.8\% \\
width & self & lead & 64.6\% & 0.003 & 55.9\% \\
width & self & fusion & 64.3\% & 0.003 & 56.1\% \\
6 obstacles & level & lead & 28.6\% & 0.003 & 20.1\% \\
6 obstacles & level & fusion & 29.1\% & 0.003 & 20.2\% \\
6 obstacles & self & lead & 45.4\% & 0.003 & 21.0\% \\
6 obstacles & self & fusion & 45.3\% & 0.003 & 20.8\% \\
\bottomrule
\end{tabular}
}
\end{adjustbox}

\end{table}

\begin{table}[H]
\centering
\caption{Ordinal model and confidence of the generic logistic regression (height).}
\label{tab:ordinal}
\begin{adjustbox}{max width=\textwidth}
{\scriptsize
\begin{tabular}{llrrrrr}
\toprule
\textbf{model} & \textbf{features} & \textbf{bal. acc} & \textbf{95\% CI} & \textbf{H1$<$$\rightarrow$H3 errors} & \textbf{MAE of class (\% LL)} & \textbf{MAE of expected height (\% LL)} \\
\midrule
nominal & lead only & 62.1\% & 55--70 & 2 & 3.9 & 4.2 \\
ordinal & lead only & 63.3\% & 55--71 & 4 & 3.8 & 4.1 \\
nominal & all & 53.7\% & 45--63 & 4 & 4.8 & 4.9 \\
ordinal & all & 54.1\% & 45--64 & 3 & 4.7 & 4.7 \\
nominal & lead foot peak height (1 feature) & 63.6\% & 52--76 & 3 & 3.8 & 4.6 \\
ordinal & lead foot peak height (1 feature) & 63.1\% & 52--75 & 2 & 3.8 & 4.7 \\
\bottomrule
\end{tabular}
}
\end{adjustbox}

\end{table}

\begin{table}[H]
\centering
\caption{Personal calibration with $k$ labelled crossings per height from the new person (logistic regression; $w$ = weight of the
person's crossings).}
\label{tab:personalcal}
\begin{adjustbox}{max width=\textwidth}
{\footnotesize
\begin{tabular}{lllrrrr}
\toprule
\textbf{norm} & \textbf{limbs} & \textbf{w} & \textbf{k = 0} & \textbf{k = 1} & \textbf{k = 2} & \textbf{k = 3} \\
\midrule
level & lead & 1 & 62.1\% & 67.6\% (65--70) & 70.8\% (69--74) & 73.6\% (71--75) \\
level & lead & 10 & 62.1\% & 72.3\% (68--79) & 76.9\% (73--80) & 79.6\% (75--83) \\
level & fusion & 1 & 56.1\% & 63.8\% (61--66) & 69.9\% (67--72) & 74.8\% (72--78) \\
level & fusion & 10 & 56.1\% & 74.4\% (71--78) & 79.3\% (76--82) & 82.2\% (80--85) \\
self & lead & 1 & 79.6\% & 81.0\% (80--82) & 81.7\% (79--84) & 82.4\% (80--86) \\
self & lead & 10 & 79.6\% & 81.6\% (80--83) & 82.8\% (81--85) & 82.9\% (80--85) \\
self & fusion & 1 & 80.8\% & 82.5\% (81--84) & 83.6\% (82--86) & 84.2\% (81--87) \\
self & fusion & 10 & 80.8\% & 83.0\% (81--86) & 85.3\% (83--87) & 85.9\% (83--88) \\
\bottomrule
\end{tabular}
}
\end{adjustbox}

\end{table}

\begin{table}[H]
\centering
\caption{Person scaling learned during use (causal): accuracy of the crossings after the first 5, 10, 15 or 20.}
\label{tab:causal}
\begin{adjustbox}{max width=\textwidth}
{\footnotesize
\begin{tabular}{llrrrr}
\toprule
\textbf{model} & \textbf{limbs} & \textbf{after 5 crossings} & \textbf{after 10} & \textbf{after 15} & \textbf{after 20} \\
\midrule
logreg C=0.1 & lead & 66.9\% (n=210) & 70.6\% (n=160) & 74.7\% (n=110) & 72.4\% (n=62) \\
logreg C=0.1 & fusion & 70.7\% (n=210) & 75.0\% (n=160) & 82.9\% (n=110) & 84.2\% (n=62) \\
forest & lead & 68.8\% (n=210) & 72.8\% (n=160) & 77.5\% (n=110) & 81.2\% (n=62) \\
forest & fusion & 69.8\% (n=210) & 74.9\% (n=160) & 81.4\% (n=110) & 88.3\% (n=62) \\
\bottomrule
\end{tabular}
}
\end{adjustbox}

\end{table}

\begin{table}[H]
\centering
\caption{Nested selection, per left-out participant: setting chosen on the training participants and its accuracy.}
\label{tab:nestedpp}
\begin{adjustbox}{max width=\textwidth}
{\footnotesize
\begin{tabular}{llrr}
\toprule
\textbf{unseen participant} & \textbf{chosen setting} & \textbf{inner accuracy} & \textbf{accuracy on the unseen participant} \\
\midrule
101 & self, forest, fusion & 90.1\% & 94\% \\
102 & self, forest, fusion & 92.1\% & 83\% \\
103 & self, forest, fusion & 88.8\% & 97\% \\
104 & self, forest, fusion & 88.9\% & 96\% \\
105 & self, forest, fusion & 88.5\% & 100\% \\
106 & self, forest, fusion & 90.0\% & 93\% \\
107 & self, forest, fusion & 91.0\% & 85\% \\
108 & self, forest, trail & 90.7\% & 62\% \\
109 & self, forest, fusion & 91.0\% & 88\% \\
110 & self, forest, fusion & 89.0\% & 96\% \\
\bottomrule
\end{tabular}
}
\end{adjustbox}

\end{table}

\begin{table}[H]
\centering
\caption{Within-person ranking: a generic model's height score ranks one person's crossings, which are split into thirds (no labels;
assumes the person crossed a mix of heights).}
\label{tab:ranking}
\begin{adjustbox}{max width=\textwidth}
{\footnotesize
\begin{tabular}{llrr}
\toprule
\textbf{features} & \textbf{limb} & \textbf{generic model} & \textbf{+ ranking into thirds} \\
\midrule
summary & lead & 62.1\% & 83.6\% (79--88; worst 67\%) \\
summary & fusion & 56.1\% & 87.0\% (84--90; worst 80\%) \\
shape & lead & 72.8\% & 85.9\% (83--89; worst 79\%) \\
shape & fusion & 63.7\% & 87.0\% (82--91; worst 72\%) \\
summary + shape & lead & 64.5\% & 83.6\% (80--87; worst 72\%) \\
summary + shape & fusion & 56.2\% & 86.2\% (82--90; worst 72\%) \\
\bottomrule
\end{tabular}
}
\end{adjustbox}

\end{table}

\begin{table}[H]
\centering
\caption{Self-training on the new person (pseudo-labels from confident predictions).}
\label{tab:selftrain}
\begin{adjustbox}{max width=\textwidth}
{\footnotesize
\begin{tabular}{lrrrrr}
\toprule
\textbf{features (lead, level)} & \textbf{no self-training} & \textbf{th 0.5, 1 round} & \textbf{th 0.5, 3 rounds} & \textbf{th 0.7, 1 round} & \textbf{th 0.7, 3 rounds} \\
\midrule
summary & 62.1\% & 62.1\% (55--70; worst 42\%) & 61.7\% (54--70; worst 42\%) & 56.7\% (52--62; worst 42\%) & 54.0\% (48--60; worst 35\%) \\
shape & 72.8\% & 65.6\% (61--71; worst 54\%) & 64.2\% (62--66; worst 59\%) & 67.1\% (62--72; worst 54\%) & 64.4\% (60--69; worst 54\%) \\
\bottomrule
\end{tabular}
}
\end{adjustbox}

\end{table}

\begin{table}[H]
\centering
\caption{Personal calibration with the clearance-shape features ($k$ labelled crossings per height).}
\label{tab:fewshot}
\begin{adjustbox}{max width=\textwidth}
{\footnotesize
\begin{tabular}{lrrrr}
\toprule
\textbf{features (lead)} & \textbf{k = 0} & \textbf{k = 1} & \textbf{k = 2} & \textbf{k = 3} \\
\midrule
summary & 62.1\% & 71.8\% (68--76) & 76.8\% (74--79) & 79.7\% (77--82) \\
shape & 72.8\% & 80.6\% (77--84) & 83.9\% (82--86) & 84.4\% (82--87) \\
summary + shape & 64.5\% & 73.6\% (71--77) & 78.5\% (74--81) & 80.8\% (78--84) \\
\bottomrule
\end{tabular}
}
\end{adjustbox}

\end{table}

\begin{table}[H]
\centering
\caption{Dose-response per participant: change per cm of obstacle (median and range over participants).}
\label{tab:dose}
\begin{adjustbox}{max width=\textwidth}
{\footnotesize
\begin{tabular}{p{5.4cm}rrrr}
\toprule
\textbf{variable} & \textbf{median slope per cm} & \textbf{range over participants} & \textbf{participants with slope $>$ 0} & \textbf{median $R^2$} \\
\midrule
lead foot lift (Dot, mm) & 10.5 & 8.5 to 11.2 & 10 / 10 & 0.87 \\
trail foot lift (Dot, mm) & 11.3 & 6.0 to 13.8 & 10 / 10 & 0.85 \\
lead min clearance (camera, mm) & $-$1.4 & $-$3.1 to 2.4 & 3 / 10 & 0.12 \\
trail min clearance (camera, mm) & $-$1.1 & $-$10.3 to 1.4 & 3 / 10 & 0.05 \\
lead thigh flexion (deg) & 1.2 & 0.9 to 1.7 & 10 / 10 & 0.72 \\
trail shank backward tilt (deg) & $-$1.1 & $-$1.7 to $-$0.2 & 0 / 10 & 0.73 \\
\bottomrule
\end{tabular}
}
\end{adjustbox}

\end{table}

\begin{table}[H]
\centering
\caption{Adaptation over the five repeats of an obstacle (deviation from the participant $\times$ obstacle mean, 95\,\% CI).}
\label{tab:adapt}
\begin{adjustbox}{max width=\textwidth}
{\scriptsize
\begin{tabular}{p{4.6cm}rrrrrr}
\toprule
\textbf{variable} & \textbf{repeat 1} & \textbf{2} & \textbf{3} & \textbf{4} & \textbf{5} & \textbf{5 minus 1} \\
\midrule
lead foot lift (Dot, mm) & +11.7 (+4.7 to +19.9) & $-$8.9 ($-$17.5 to +1.1) & +1.7 ($-$4.9 to +9.4) & $-$6.3 ($-$19.0 to +4.5) & +5.0 (+0.5 to +9.7) & $-$6.7 \\
trail foot lift (Dot, mm) & +2.1 ($-$15.0 to +16.6) & +3.0 ($-$9.7 to +15.5) & $-$8.1 ($-$14.4 to $-$1.9) & $-$3.6 ($-$13.0 to +7.9) & $-$1.7 ($-$14.5 to +8.7) & $-$3.8 \\
lead min clearance (camera, mm) & +8.6 ($-$1.3 to +22.3) & $-$14.0 ($-$30.6 to $-$2.6) & +5.5 ($-$3.2 to +17.0) & $-$10.3 ($-$23.2 to $-$0.6) & +9.4 (+2.2 to +16.0) & +0.8 \\
trail min clearance (camera, mm) & $-$9.5 ($-$26.0 to +7.9) & +1.5 ($-$13.4 to +18.7) & +3.5 ($-$15.6 to +25.5) & +7.7 ($-$8.9 to +26.9) & $-$6.2 ($-$22.4 to +7.4) & +3.4 \\
lead thigh flexion (deg) & $-$0.2 ($-$2.1 to +2.0) & $-$1.9 ($-$3.7 to $-$0.7) & +1.5 (+0.1 to +2.9) & $-$0.3 ($-$2.5 to +1.3) & +0.9 ($-$0.1 to +1.9) & +1.1 \\
trail shank backward tilt (deg) & $-$1.7 ($-$3.8 to +0.4) & $-$0.3 ($-$2.7 to +1.8) & +1.2 (+0.2 to +2.3) & +1.0 ($-$0.5 to +2.3) & +0.4 ($-$1.5 to +2.2) & +2.1 \\
\bottomrule
\end{tabular}
}
\end{adjustbox}

\end{table}

\begin{table}[H]
\centering
\caption{Features most driven by the obstacle height within person, and how much the person explains for them.}
\label{tab:persontop}
\begin{adjustbox}{max width=\textwidth}
{\footnotesize
\begin{tabular}{lrr}
\toprule
\textbf{feature} & \textbf{$\eta^2$ height within person} & \textbf{$\eta^2$ person} \\
\midrule
\texttt{trail\_\allowbreak{}zhc\_\allowbreak{}peak\_\allowbreak{}foot\_\allowbreak{}h} & 0.77 & 0.37 \\
\texttt{lead\_\allowbreak{}zhc\_\allowbreak{}peak\_\allowbreak{}foot\_\allowbreak{}h} & 0.76 & 0.36 \\
\texttt{trail\_\allowbreak{}zhc\_\allowbreak{}h\_\allowbreak{}25} & 0.75 & 0.38 \\
\texttt{lead\_\allowbreak{}thigh\_\allowbreak{}max} & 0.74 & 0.36 \\
\texttt{trail\_\allowbreak{}shank\_\allowbreak{}min} & 0.70 & 0.55 \\
\texttt{lead\_\allowbreak{}thigh\_\allowbreak{}max\_\allowbreak{}corr} & 0.67 & 0.33 \\
\texttt{trail\_\allowbreak{}zhc\_\allowbreak{}h\_\allowbreak{}toeoff} & 0.67 & 0.36 \\
\texttt{lead\_\allowbreak{}zhc\_\allowbreak{}h\_\allowbreak{}50} & 0.62 & 0.46 \\
\texttt{trail\_\allowbreak{}shank\_\allowbreak{}min\_\allowbreak{}corr} & 0.58 & 0.49 \\
\texttt{trail\_\allowbreak{}shank\_\allowbreak{}rom} & 0.57 & 0.44 \\
\bottomrule
\end{tabular}
}
\end{adjustbox}

\end{table}

\begin{table}[H]
\centering
\caption{Which leg leads: per participant, dominant leg, side that led, share led with the dominant leg, by height.}
\label{tab:leads}
\begin{adjustbox}{max width=\textwidth}
{\footnotesize
\begin{tabular}{llrrrrr}
\toprule
\textbf{participant} & \textbf{dominant} & \textbf{led with R / L} & \textbf{led with dominant} & \textbf{H1} & \textbf{H2} & \textbf{H3} \\
\midrule
101 & L & 19 / 13 & 41\% & 36\% & 30\% & 55\% \\
102 & R & 12 / 18 & 40\% & 40\% & 40\% & 40\% \\
103 & R & 15 / 15 & 50\% & 30\% & 60\% & 60\% \\
104 & R & 30 / 0 & 100\% & 100\% & 100\% & 100\% \\
105 & R & 14 / 16 & 47\% & 60\% & 40\% & 40\% \\
106 & L & 22 / 8 & 27\% & 40\% & 0\% & 40\% \\
107 & R & 4 / 26 & 13\% & 0\% & 10\% & 30\% \\
108 & R & 13 / 17 & 43\% & 70\% & 20\% & 40\% \\
109 & R & 13 / 18 & 42\% & 40\% & 40\% & 45\% \\
110 & R & 15 / 15 & 50\% & 30\% & 50\% & 70\% \\
\textbf{all} & | 157 / 146 & 45\% & 45\% & 39\% & 52\% &  \\
\bottomrule
\end{tabular}
}
\end{adjustbox}

\end{table}

\begin{table}[H]
\centering
\caption{Camera-free pipeline: all-correct crossings per participant.}
\label{tab:streampp}
\begin{adjustbox}{max width=\textwidth}
{\footnotesize
\begin{tabular}{lrrrrrrrrrr}
\toprule
\textbf{participant} & \textbf{101} & \textbf{102} & \textbf{103} & \textbf{104} & \textbf{105} & \textbf{106} & \textbf{107} & \textbf{108} & \textbf{109} & \textbf{110} \\
\midrule
all correct, crossing-scaled & 75\% & 83\% & 93\% & 93\% & 96\% & 90\% & 93\% & 69\% & 93\% & 80\% \\
\bottomrule
\end{tabular}
}
\end{adjustbox}

\end{table}

\begin{table}[H]
\centering
\caption{Camera-free pipeline, learning during use: each crossing scaled with the earlier detected crossings only.}
\label{tab:streamcausal}
\begin{adjustbox}{max width=\textwidth}
{\footnotesize
\begin{tabular}{lrrr}
\toprule
\textbf{warm-up (earlier detections)} & \textbf{crossings scored} & \textbf{height right (when found)} & \textbf{all correct} \\
\midrule
5 & 244 & 69.1\% & 65.2\% \\
10 & 196 & 74.7\% & 69.4\% \\
15 & 149 & 85.2\% & 77.2\% \\
20 & 102 & 87.5\% & 75.5\% \\
\bottomrule
\end{tabular}
}
\end{adjustbox}

\end{table}
